% Originally: content/week-04.tex, \subsection{What does imply differentiability} --- promoted

% Supplement: Sascha Brack/Class Notes/Week_04_Notes_Monday_Updated.pdf, pp. 3--5
\section{What \emph{does} imply differentiability}
\label{sec:sufficient_condition_differentiability}

\begin{ainote}
Corsin's notes give the definition and the warnings, but never state a \emph{sufficient}
condition --- so far we have two counterexamples and no positive test. Sascha Brack's notes
(Lecture 8) state one as Theorem 2.9 and organise the whole picture into a chain of
implications; both are reproduced below, since without them the practical question
\qt{how do I actually check differentiability?} is left unanswered.
\end{ainote}

\begin{theorem}[Continuous partial derivatives suffice]
\label{thm:continuous_partials_imply_differentiable}
Let $U \subseteq \mathbb{R}^n$ be open and $F : U \to \mathbb{R}^m$. If all partial derivatives
$\partial_iF$ exist on $U$ \textbf{and are continuous} at every $x \in U$, then $F$ is
differentiable at every $x \in U$, and
\[DF_x = \big(\partial_1F(x),\ \dots,\ \partial_nF(x)\big).\]
Such an $F$ is called \newterm{continuously differentiable}, written $F \in C^1(U)$.
\end{theorem}

This is the theorem that makes the subject workable: in practice one almost never checks the
definition directly. Compute the partials, observe that they are built from continuous pieces,
and differentiability follows --- which is exactly why the course can write $f \in C^1$, $C^2$,
$C^k$ so freely from here onward.

Together with the two counterexamples, the whole landscape is a single chain, each link
one-directional:

% Sonnet 5 (Medium) --- FIG-W04-02, after the diagram on Sascha Brack's p. 3.
\begin{center}
\begin{tikzpicture}[>=Latex, node distance=0pt]
  \tikzset{
    box/.style={draw, rounded corners=2pt, align=center, font=\small,
                inner sep=4pt, minimum height=8mm},
    imp/.style={->, very thick, green!45!black},
    nimp/.style={->, thick, red!75!black, dashed}}

  \node[box, fill=green!8] (c1) at (0,4.5) {$F \in C^1$ near $x_0$\\[-2pt]
    {\scriptsize (all $\partial_iF$ exist \emph{and} are continuous)}};
  \node[box, fill=blue!8] (df) at (0,2.9) {$F$ differentiable at $x_0$};
  \node[box] (dir) at (0,1.4) {all directional derivatives $\partial_vF(x_0)$ exist};
  \node[box] (par) at (0,-0.1) {all partial derivatives $\partial_iF(x_0)$ exist};

  \draw[imp] (c1) -- (df);
  \draw[imp] (df) -- (dir);
  \draw[imp] (dir) -- (par);

  \node[font=\scriptsize, green!45!black, right] at (0.15,3.7)
    {\cref{thm:continuous_partials_imply_differentiable}};
  \node[font=\scriptsize, green!45!black, right] at (0.15,2.15) {formula $(*)$};
  \node[font=\scriptsize, green!45!black, right] at (0.15,0.65) {take $v = e_i$};

  % The reverse arrows, all false. Each carries its counterexample as a midway node on the
  % arc itself, filled white so it interrupts the dashed line cleanly. Previously the three
  % labels were free-standing nodes anchored at x = -2.6 while the arcs bulged out past
  % x = -3, so every arrow was drawn straight through its own label.
  \tikzset{noLbl/.style={midway, font=\scriptsize, red!75!black, align=right,
                         fill=white, inner sep=1.5pt}}
  \draw[nimp] (df.west)  to[out=160, in=200, looseness=1.7]
    node[noLbl] {\textbf{no:} \cref{ex:differentiable_not_c1}} (c1.west);
  \draw[nimp] (dir.west) to[out=160, in=200, looseness=1.7]
    node[noLbl] {\textbf{no:} \cref{ex:all_directional_derivatives_exist}} (df.west);
  \draw[nimp] (par.west) to[out=160, in=200, looseness=1.7]
    node[noLbl] {\textbf{no:} \cref{ex:partial_derivatives_not_continuous}} (dir.west);
\end{tikzpicture}
\end{center}

Every downward arrow holds, and \textbf{not one of them reverses} --- each failure is witnessed
by an example already in this chapter, so the diagram doubles as a map of where those
counterexamples belong:
\begin{itemize}
  \item \textbf{Differentiable $\not\Rightarrow$ $C^1$:} \cref{ex:differentiable_not_c1}, whose
    partial derivatives oscillate without limit at the origin even though the differential
    exists there.
  \item \textbf{All directional derivatives $\not\Rightarrow$ differentiable:}
    \cref{ex:all_directional_derivatives_exist}, where every ray gives a slope but
    $v \mapsto \partial_vf(0,0)$ is not linear.
  \item \textbf{Partial derivatives $\not\Rightarrow$ all directional derivatives:}
    \cref{ex:partial_derivatives_not_continuous} again. For $f = \tfrac{xy}{x^2+y^2}$ and a
    direction $v = (a,b)$ with $ab \neq 0$,
    \[\frac{f(ta,tb)}{t} = \frac{1}{t}\cdot\frac{t^2ab}{t^2(a^2+b^2)}
      = \frac{ab}{t\,(a^2+b^2)} \ \xrightarrow{\ t \to 0\ }\ \pm\infty,\]
    so $\partial_vf(0,0)$ exists only along the two axes --- the two directions the partial
    derivatives happen to look at.
\end{itemize}
The three counterexamples are not variations on one theme; they fail at three different
heights, which is precisely why the chain has three separate links.

% Supplement: Sascha Brack/Class Notes/Week_04_Notes_Monday_Updated.pdf, p. 5
\begin{remark}[A recipe for checking differentiability at a point]
\label{rem:differentiability_recipe}
Sascha's notes distil the above into three steps, worth following in this order.
\begin{enumerate}[label=\textbf{\arabic*.}]
  \item \textbf{Compute the partial derivatives $\partial_iF(x_0)$.} If one of them fails to
    exist, $F$ is not differentiable and you are done. If they all exist \emph{and are
    continuous near $x_0$}, $F$ \emph{is} differentiable by
    \Cref{thm:continuous_partials_imply_differentiable} --- also done. Most exercises end here.
  \item \textbf{Otherwise, guess the candidate.} If $F$ is differentiable at all, the differential
    can only be $L := \big(\partial_1F(x_0),\dots,\partial_nF(x_0)\big)$, since the partials are
    forced (\cref{sec:jacobi_matrix}).
  \item \textbf{Verify the candidate against the definition:}
    \[\lim_{h\to0}\frac{\big|F(x_0+h)-F(x_0)-L(h)\big|}{\lVert h\rVert} \overset{?}{=} 0.\]
    In $\mathbb{R}^2$, substituting $h = (r\cos\theta,\ r\sin\theta)$ and checking that the
    result tends to $0$ \emph{uniformly in $\theta$} is usually the quickest route --- the same
    polar trick used in \cref{ex:all_directional_derivatives_exist} and, later, in
    \Cref{q:quiz_q4}.
\end{enumerate}
Step 1 is the one students skip; it settles the overwhelming majority of cases without any
$\varepsilon$--$\delta$ work at all.
\end{remark}

% Supplement: Sascha Brack/Class Notes/Week_03_Notes_Friday.pdf, p. 7
\begin{example}[Checking differentiability at the origin]
\label{ex:differentiability_recipe_application}
Let $f(x,y) := \frac{x^2 y}{x^2+y^2}$ for $(x,y) \neq (0,0)$ and $f(0,0) := 0$. We check if $f$ is differentiable at the origin.
\begin{enumerate}[label=\textbf{\arabic*.}]
  \item \textbf{Partial derivatives:}
    \[\partial_1 f(0,0) = \lim_{t\to 0}\frac{f(t,0)-f(0,0)}{t} = \lim_{t\to 0}\frac{0}{t} = 0, \qquad \partial_2 f(0,0) = \lim_{t\to 0}\frac{f(0,t)-f(0,0)}{t} = 0.\]
    They exist, but to use \cref{thm:continuous_partials_imply_differentiable} we would have to compute them everywhere and show they are continuous at $(0,0)$, which is tedious and (as we will see) false. We proceed to step 2.
  \item \textbf{Guess the candidate:} The only possible differential is $L(h_1,h_2) = \partial_1 f(0,0)h_1 + \partial_2 f(0,0)h_2 = 0$.
  \item \textbf{Verify against the definition:} We check whether the following limit is zero:
    \[\lim_{h\to 0} \frac{|f(h_1,h_2) - f(0,0) - L(h)|}{\lVert h\rVert} = \lim_{h\to 0} \frac{\frac{h_1^2 h_2}{h_1^2+h_2^2}}{\sqrt{h_1^2+h_2^2}} = \lim_{h\to 0} \frac{h_1^2 h_2}{(h_1^2+h_2^2)^{3/2}}.\]
    Passing to polar coordinates $h = (r\cos\theta, r\sin\theta)$, the expression becomes
    \[\frac{r^3\cos^2\theta\sin\theta}{r^3} = \cos^2\theta\sin\theta.\]
    This depends on the angle of approach $\theta$ (for instance, it is $0$ along the axes but $\frac{1}{2\sqrt{2}} \neq 0$ along $h_1=h_2$), so the limit as $r \to 0$ is not uniformly $0$. Hence $f$ is \textbf{not differentiable} at $(0,0)$.
\end{enumerate}
\end{example}

% Supplement: Sascha Brack/Class Notes/Week_04_Notes_Friday_Updated.pdf, pp. 4--5
\begin{example}[A function that is differentiable at the origin]
\label{ex:differentiability_recipe_positive}
Let $f(x,y) := \frac{x^2 y^2}{\sqrt{x^2+y^2}}$ for $(x,y) \neq (0,0)$ and $f(0,0) := 0$. We check differentiability at the origin following the recipe.
\begin{enumerate}[label=\textbf{\arabic*.}]
  \item \textbf{Partial derivatives:} For $(x,y) \neq (0,0)$, the quotient and chain rules give
    \[\partial_1 f(x,y) = xy^2\frac{x^2+2y^2}{(x^2+y^2)^{3/2}}, \qquad \partial_2 f(x,y) = x^2y\frac{2x^2+y^2}{(x^2+y^2)^{3/2}}.\]
    At the origin, computing directly from the definition yields $\partial_1 f(0,0) = 0$ and $\partial_2 f(0,0) = 0$.
    To see if they are continuous at $(0,0)$, we substitute $x = r\cos\theta$ and $y = r\sin\theta$:
    \[\partial_1 f(r\cos\theta, r\sin\theta) = \frac{r\cos\theta\,r^2\sin^2\theta \cdot r^2(1+\sin^2\theta)}{r^3} = r^2\cos\theta\sin^2\theta(1+\sin^2\theta).\]
    As $r \to 0$, this tends to $0$ uniformly in $\theta$, so $\partial_1 f$ is continuous at $(0,0)$. The same holds for $\partial_2 f$. By \cref{thm:continuous_partials_imply_differentiable}, $f$ is differentiable at the origin.
  \item \textbf{Guess the candidate:} We already know $f$ is differentiable and $L = (0,0)$. But for practice, let's verify it with the definition directly.
  \item \textbf{Verify against the definition:} We check the limit
    \[\lim_{h\to 0} \frac{|f(h_1,h_2) - f(0,0) - L(h)|}{\lVert h\rVert} = \lim_{h\to 0} \frac{\frac{h_1^2 h_2^2}{\sqrt{h_1^2+h_2^2}}}{\sqrt{h_1^2+h_2^2}} = \lim_{h\to 0} \frac{h_1^2 h_2^2}{h_1^2+h_2^2}.\]
    In polar coordinates, this is $\frac{r^4\cos^2\theta\sin^2\theta}{r^2} = r^2\cos^2\theta\sin^2\theta$. As $r \to 0$, the limit is clearly $0$ uniformly in $\theta$, confirming $f$ is differentiable with $Df_0 = 0$.
\end{enumerate}
\end{example}
